
\documentclass{article}
\usepackage{colortbl}
\usepackage{makecell}
\usepackage{multirow}
\usepackage{supertabular}

\begin{document}

\newcounter{utterance}

\centering \large Interaction Transcript for game `codenames', experiment `opponent\_easy', episode 2 with glm{-}5.2.
\vspace{24pt}

{ \footnotesize  \setcounter{utterance}{1}
\setlength{\tabcolsep}{0pt}
\begin{supertabular}{c@{$\;$}|p{.15\linewidth}@{}p{.15\linewidth}p{.15\linewidth}p{.15\linewidth}p{.15\linewidth}p{.15\linewidth}}
    \# & $\;$A & \multicolumn{4}{c}{Game Master} & $\;\:$B\\
    \hline

    \theutterance \stepcounter{utterance}  
    & & & \multicolumn{2}{p{0.3\linewidth}}{
        \cellcolor[rgb]{0.9,0.9,0.9}{
            \makecell[{{p{\linewidth}}}]{
                \texttt{\tiny{[GM$|$GM]}}
                {'hidden': {'team': ['taniec', 'port', 'indyk', 'ryba', 'miedź', 'pralka', 'francja', 'laboratorium', 'bugle'], 'innocent': ['helikopter', 'ośmiornica', 'australia', 'twarz', 'czekolada', 'zegarek', 'kłoda'], 'opponent': ['mucha', 'dyrygent', 'talerz', 'oliwka', 'dobrze', 'podanie', 'spadochron', 'cykl'], 'assassin': ['choroba']}, 'revealed': {'team': {'team': [], 'innocent': [], 'opponent': [], 'assassin': []}, 'opponent': {'team': [], 'innocent': [], 'opponent': [], 'assassin': []}}}
            }
        }
    }
    & & \\ \\

    \theutterance \stepcounter{utterance}  
    & & \multicolumn{4}{p{0.6\linewidth}}{
        \cellcolor[rgb]{0.9,0.9,0.9}{
            \makecell[{{p{\linewidth}}}]{
                \texttt{\tiny{[P1$\langle$GM]}}
                \texttt{Podaj jedną wskazówkę składającą się z pojedynczego słowa, która odnosi się do niektórych słów twojej drużyny z poniższej listy. Możesz wybrać, do ilu i do jakich słów odnosi się twoja wskazówka.} \\
\texttt{Celem jest, aby twój partner odgadł wszystkie słowa twojej drużyny na podstawie podanych wskazówek, zanim drużyna przeciwna odkryje wszystkie swoje. Jeśli twój partner odgadnie słowo 'przeciwnika' lub 'rozpraszające', jego tura zgadywania kończy się i drużyna przeciwna wykonuje turę. Jeśli twój partner odgadnie słowo 'zabójcze', gra kończy się i twoja drużyna natychmiast przegrywa.} \\
\texttt{Im więcej słów odnosi się do twojej wskazówki, tym szybciej możesz wygrać grę, ale tym trudniejsze mogą być skojarzenia do odgadnięcia dla twojego partnera, więc wybieraj słowo wskazówki mądrze. Słowo wskazówki musi być semantycznie powiązane ze słowami docelowymi, nie może być jednym ze słów na listach ani zawierać ich części.} \\
\\ 
\texttt{Zawsze podawaj swoją wskazówkę składającą się z pojedynczego słowa i oddzieloną przecinkami listę powiązanych słów docelowych w następującym formacie i staraj się, aby twoje odpowiedzi były jak najkrótsze, nigdy nie dołączaj żadnego innego tekstu niż jest to wymagane w tym formacie:} \\
\\ 
\texttt{WSKAZÓWKA: <SŁOWO>} \\
\texttt{CELE: <CELE>} \\
\\ 
\texttt{Słowa twojej drużyny to: taniec, port, indyk, ryba, miedź, pralka, francja, laboratorium, bugle.} \\
\texttt{Słowa drużyny przeciwnika to: mucha, dyrygent, talerz, oliwka, dobrze, podanie, spadochron, cykl.} \\
\texttt{Słowa rozpraszające to: helikopter, ośmiornica, australia, twarz, czekolada, zegarek, kłoda.} \\
\texttt{Słowa zabójcze to: choroba.} \\
            }
        }
    }
    & \\ \\

    \theutterance \stepcounter{utterance}  
    & \multicolumn{4}{p{0.6\linewidth}}{
        \cellcolor[rgb]{0.9,0.9,0.9}{
            \makecell[{{p{\linewidth}}}]{
                \texttt{\tiny{[P1$\rangle$GM]}}
                \texttt{WSKAZÓWKA: sprzęt} \\
\texttt{CELE: pralka, laboratorium} \\
            }
        }
    }
    & & \\ \\

    \theutterance \stepcounter{utterance}  
    & & & \multicolumn{2}{p{0.3\linewidth}}{
        \cellcolor[rgb]{0.9,0.9,0.9}{
            \makecell[{{p{\linewidth}}}]{
                \texttt{\tiny{[GM$|$GM]}}
                \texttt{sprzęt} \\
            }
        }
    }
    & & \\ \\

    \theutterance \stepcounter{utterance}  
    & & & \multicolumn{2}{p{0.3\linewidth}}{
        \cellcolor[rgb]{0.9,0.9,0.9}{
            \makecell[{{p{\linewidth}}}]{
                \texttt{\tiny{[GM$|$GM]}}
                ['pralka', 'laboratorium']
            }
        }
    }
    & & \\ \\

    \theutterance \stepcounter{utterance}  
    & & & \multicolumn{2}{p{0.3\linewidth}}{
        \cellcolor[rgb]{0.9,0.9,0.9}{
            \makecell[{{p{\linewidth}}}]{
                \texttt{\tiny{[GM$|$GM]}}
                {'word': 'pralka', 'assignment': 'team'}
            }
        }
    }
    & & \\ \\

    \theutterance \stepcounter{utterance}  
    & & & \multicolumn{2}{p{0.3\linewidth}}{
        \cellcolor[rgb]{0.9,0.9,0.9}{
            \makecell[{{p{\linewidth}}}]{
                \texttt{\tiny{[GM$|$GM]}}
                {'word': 'laboratorium', 'assignment': 'team'}
            }
        }
    }
    & & \\ \\

    \theutterance \stepcounter{utterance}  
    & & \multicolumn{4}{p{0.6\linewidth}}{
        \cellcolor[rgb]{0.9,0.9,0.9}{
            \makecell[{{p{\linewidth}}}]{
                \texttt{\tiny{[GM$\rangle$P2]}}
                \texttt{Podaj oddzieloną przecinkami listę maksymalnie 2 słów z poniższej listy, które najlepiej odnoszą się lub są najbardziej związane ze słowem 'sprzęt'. Zawsze rozpoczynaj swoją listę odgadnięć od ZGADNIĘCIE:  i nie dołączaj żadnego innego tekstu w swojej odpowiedzi.} \\
\\ 
\texttt{oliwka, mucha, pralka, australia, indyk, bugle, taniec, ośmiornica, ryba, miedź, laboratorium, choroba, podanie, cykl, helikopter, zegarek, czekolada, port, spadochron, francja, dyrygent, dobrze, talerz, kłoda, twarz} \\
            }
        }
    }
    & \\ \\

    \theutterance \stepcounter{utterance}  
    & & & \multicolumn{4}{p{0.6\linewidth}}{
        \cellcolor[rgb]{0.9,0.9,0.9}{
            \makecell[{{p{\linewidth}}}]{
                \texttt{\tiny{[GM$\langle$P2]}}
                \texttt{ZGADNIĘCIE: pralka, helikopter} \\
            }
        }
    }
     \\ \\

    \theutterance \stepcounter{utterance}  
    & & & \multicolumn{2}{p{0.3\linewidth}}{
        \cellcolor[rgb]{0.9,0.9,0.9}{
            \makecell[{{p{\linewidth}}}]{
                \texttt{\tiny{[GM$|$GM]}}
                ['pralka', 'helikopter']
            }
        }
    }
    & & \\ \\

    \theutterance \stepcounter{utterance}  
    & & & \multicolumn{2}{p{0.3\linewidth}}{
        \cellcolor[rgb]{0.9,0.9,0.9}{
            \makecell[{{p{\linewidth}}}]{
                \texttt{\tiny{[GM$|$GM]}}
                {'word': 'pralka', 'assignment': 'team'}
            }
        }
    }
    & & \\ \\

    \theutterance \stepcounter{utterance}  
    & & & \multicolumn{2}{p{0.3\linewidth}}{
        \cellcolor[rgb]{0.9,0.9,0.9}{
            \makecell[{{p{\linewidth}}}]{
                \texttt{\tiny{[GM$|$GM]}}
                {'word': 'pralka', 'assignment': 'team'}
            }
        }
    }
    & & \\ \\

    \theutterance \stepcounter{utterance}  
    & & & \multicolumn{2}{p{0.3\linewidth}}{
        \cellcolor[rgb]{0.9,0.9,0.9}{
            \makecell[{{p{\linewidth}}}]{
                \texttt{\tiny{[GM$|$GM]}}
                {'word': 'helikopter', 'assignment': 'innocent'}
            }
        }
    }
    & & \\ \\

    \theutterance \stepcounter{utterance}  
    & & & \multicolumn{2}{p{0.3\linewidth}}{
        \cellcolor[rgb]{0.9,0.9,0.9}{
            \makecell[{{p{\linewidth}}}]{
                \texttt{\tiny{[GM$|$GM]}}
                \texttt{helikopter} \\
            }
        }
    }
    & & \\ \\

    \theutterance \stepcounter{utterance}  
    & & & \multicolumn{2}{p{0.3\linewidth}}{
        \cellcolor[rgb]{0.9,0.9,0.9}{
            \makecell[{{p{\linewidth}}}]{
                \texttt{\tiny{[GM$|$GM]}}
                {'word': 'spadochron', 'assignment': 'opponent'}
            }
        }
    }
    & & \\ \\

    \theutterance \stepcounter{utterance}  
    & & & \multicolumn{2}{p{0.3\linewidth}}{
        \cellcolor[rgb]{0.9,0.9,0.9}{
            \makecell[{{p{\linewidth}}}]{
                \texttt{\tiny{[GM$|$GM]}}
                {'hidden': {'team': ['taniec', 'port', 'indyk', 'ryba', 'miedź', 'francja', 'laboratorium', 'bugle'], 'innocent': ['ośmiornica', 'australia', 'twarz', 'czekolada', 'zegarek', 'kłoda'], 'opponent': ['mucha', 'dyrygent', 'talerz', 'oliwka', 'dobrze', 'podanie', 'cykl'], 'assassin': ['choroba']}, 'revealed': {'team': {'team': ['pralka'], 'innocent': ['helikopter'], 'opponent': [], 'assassin': []}, 'opponent': {'team': [], 'innocent': [], 'opponent': ['spadochron'], 'assassin': []}}}
            }
        }
    }
    & & \\ \\

    \theutterance \stepcounter{utterance}  
    & & \multicolumn{4}{p{0.6\linewidth}}{
        \cellcolor[rgb]{0.9,0.9,0.9}{
            \makecell[{{p{\linewidth}}}]{
                \texttt{\tiny{[P1$\langle$GM]}}
                \texttt{Słowo pralka zostało odgadnięte prawidłowo. Słowo helikopter zostało odgadnięte, ale to słowo niewinny. Tura twojego kolegi z drużyny się tam skończyła.} \\
\texttt{Teraz podaj kolejną wskazówkę odnoszącą się do niektórych z pozostałych słów twojego zespołu oraz listę powiązanych słów docelowych. Pamiętaj, aby rozpocząć wskazówkę od 'WSKAZÓWKA: ', wstawić nową linię i rozpocząć listę słów docelowych oddzielonych przecinkami od 'CELE: '. Uwaga: niektóre słowa zostały usunięte z list w porównaniu do poprzednich żądań.} \\
\\ 
\texttt{Twoje pozostałe słowa zespołu to: taniec, port, indyk, ryba, miedź, francja, laboratorium, bugle.} \\
\texttt{Pozostałe słowa twojego przeciwnika to: mucha, dyrygent, talerz, oliwka, dobrze, podanie, spadochron, cykl.} \\
\texttt{Pozostałe słowa rozpraszające to: ośmiornica, australia, twarz, czekolada, zegarek, kłoda.} \\
\texttt{Pozostałe słowa zabójcy to: choroba.} \\
            }
        }
    }
    & \\ \\

    \theutterance \stepcounter{utterance}  
    & \multicolumn{4}{p{0.6\linewidth}}{
        \cellcolor[rgb]{0.9,0.9,0.9}{
            \makecell[{{p{\linewidth}}}]{
                \texttt{\tiny{[P1$\rangle$GM]}}
                \texttt{WSKAZÓWKA: instrument} \\
\texttt{CELE: bugle} \\
            }
        }
    }
    & & \\ \\

    \theutterance \stepcounter{utterance}  
    & & & \multicolumn{2}{p{0.3\linewidth}}{
        \cellcolor[rgb]{0.9,0.9,0.9}{
            \makecell[{{p{\linewidth}}}]{
                \texttt{\tiny{[GM$|$GM]}}
                \texttt{instrument} \\
            }
        }
    }
    & & \\ \\

    \theutterance \stepcounter{utterance}  
    & & & \multicolumn{2}{p{0.3\linewidth}}{
        \cellcolor[rgb]{0.9,0.9,0.9}{
            \makecell[{{p{\linewidth}}}]{
                \texttt{\tiny{[GM$|$GM]}}
                ['bugle']
            }
        }
    }
    & & \\ \\

    \theutterance \stepcounter{utterance}  
    & & & \multicolumn{2}{p{0.3\linewidth}}{
        \cellcolor[rgb]{0.9,0.9,0.9}{
            \makecell[{{p{\linewidth}}}]{
                \texttt{\tiny{[GM$|$GM]}}
                {'word': 'bugle', 'assignment': 'team'}
            }
        }
    }
    & & \\ \\

    \theutterance \stepcounter{utterance}  
    & & \multicolumn{4}{p{0.6\linewidth}}{
        \cellcolor[rgb]{0.9,0.9,0.9}{
            \makecell[{{p{\linewidth}}}]{
                \texttt{\tiny{[GM$\rangle$P2]}}
                \texttt{Słowo pralka zostało odgadnięte prawidłowo. Słowo helikopter zostało odgadnięte, ale to słowo niewinny. Twoja tura się tam skończyła.} \\
\texttt{Teraz podaj kolejną listę oddzieloną przecinkami, zawierającą co najmniej 1 i maksymalnie 1 słów z poniższej listy słów, które najlepiej odnoszą się lub są najbardziej związane ze słowem 'instrument'. Pamiętaj, aby rozpocząć odpowiedź od ZGADNIĘCIE: . Uwaga: niektóre słowa zostały usunięte z listy w porównaniu do poprzednich próśb.} \\
\\ 
\texttt{oliwka, mucha, australia, indyk, bugle, taniec, ośmiornica, ryba, miedź, laboratorium, choroba, podanie, cykl, zegarek, czekolada, port, francja, dyrygent, dobrze, talerz, kłoda, twarz} \\
            }
        }
    }
    & \\ \\

    \theutterance \stepcounter{utterance}  
    & & & \multicolumn{4}{p{0.6\linewidth}}{
        \cellcolor[rgb]{0.9,0.9,0.9}{
            \makecell[{{p{\linewidth}}}]{
                \texttt{\tiny{[GM$\langle$P2]}}
                \texttt{ZGADNIĘCIE: bugle} \\
            }
        }
    }
     \\ \\

    \theutterance \stepcounter{utterance}  
    & & & \multicolumn{2}{p{0.3\linewidth}}{
        \cellcolor[rgb]{0.9,0.9,0.9}{
            \makecell[{{p{\linewidth}}}]{
                \texttt{\tiny{[GM$|$GM]}}
                ['bugle']
            }
        }
    }
    & & \\ \\

    \theutterance \stepcounter{utterance}  
    & & & \multicolumn{2}{p{0.3\linewidth}}{
        \cellcolor[rgb]{0.9,0.9,0.9}{
            \makecell[{{p{\linewidth}}}]{
                \texttt{\tiny{[GM$|$GM]}}
                {'word': 'bugle', 'assignment': 'team'}
            }
        }
    }
    & & \\ \\

    \theutterance \stepcounter{utterance}  
    & & & \multicolumn{2}{p{0.3\linewidth}}{
        \cellcolor[rgb]{0.9,0.9,0.9}{
            \makecell[{{p{\linewidth}}}]{
                \texttt{\tiny{[GM$|$GM]}}
                {'word': 'bugle', 'assignment': 'team'}
            }
        }
    }
    & & \\ \\

    \theutterance \stepcounter{utterance}  
    & & & \multicolumn{2}{p{0.3\linewidth}}{
        \cellcolor[rgb]{0.9,0.9,0.9}{
            \makecell[{{p{\linewidth}}}]{
                \texttt{\tiny{[GM$|$GM]}}
                {'word': 'talerz', 'assignment': 'opponent'}
            }
        }
    }
    & & \\ \\

    \theutterance \stepcounter{utterance}  
    & & & \multicolumn{2}{p{0.3\linewidth}}{
        \cellcolor[rgb]{0.9,0.9,0.9}{
            \makecell[{{p{\linewidth}}}]{
                \texttt{\tiny{[GM$|$GM]}}
                {'hidden': {'team': ['taniec', 'port', 'indyk', 'ryba', 'miedź', 'francja', 'laboratorium'], 'innocent': ['ośmiornica', 'australia', 'twarz', 'czekolada', 'zegarek', 'kłoda'], 'opponent': ['mucha', 'dyrygent', 'oliwka', 'dobrze', 'podanie', 'cykl'], 'assassin': ['choroba']}, 'revealed': {'team': {'team': ['pralka', 'bugle'], 'innocent': ['helikopter'], 'opponent': [], 'assassin': []}, 'opponent': {'team': [], 'innocent': [], 'opponent': ['spadochron', 'talerz'], 'assassin': []}}}
            }
        }
    }
    & & \\ \\

    \theutterance \stepcounter{utterance}  
    & & \multicolumn{4}{p{0.6\linewidth}}{
        \cellcolor[rgb]{0.9,0.9,0.9}{
            \makecell[{{p{\linewidth}}}]{
                \texttt{\tiny{[P1$\langle$GM]}}
                \texttt{Słowo bugle zostało odgadnięte prawidłowo. Tura twojego kolegi z drużyny się tam skończyła.} \\
\texttt{Teraz podaj kolejną wskazówkę odnoszącą się do niektórych z pozostałych słów twojego zespołu oraz listę powiązanych słów docelowych. Pamiętaj, aby rozpocząć wskazówkę od 'WSKAZÓWKA: ', wstawić nową linię i rozpocząć listę słów docelowych oddzielonych przecinkami od 'CELE: '. Uwaga: niektóre słowa zostały usunięte z list w porównaniu do poprzednich żądań.} \\
\\ 
\texttt{Twoje pozostałe słowa zespołu to: taniec, port, indyk, ryba, miedź, francja, laboratorium.} \\
\texttt{Pozostałe słowa twojego przeciwnika to: mucha, dyrygent, talerz, oliwka, dobrze, podanie, cykl.} \\
\texttt{Pozostałe słowa rozpraszające to: ośmiornica, australia, twarz, czekolada, zegarek, kłoda.} \\
\texttt{Pozostałe słowa zabójcy to: choroba.} \\
            }
        }
    }
    & \\ \\

    \theutterance \stepcounter{utterance}  
    & \multicolumn{4}{p{0.6\linewidth}}{
        \cellcolor[rgb]{0.9,0.9,0.9}{
            \makecell[{{p{\linewidth}}}]{
                \texttt{\tiny{[P1$\rangle$GM]}}
                \texttt{WSKAZÓWKA: kraj} \\
\texttt{CELE: francja} \\
            }
        }
    }
    & & \\ \\

    \theutterance \stepcounter{utterance}  
    & & & \multicolumn{2}{p{0.3\linewidth}}{
        \cellcolor[rgb]{0.9,0.9,0.9}{
            \makecell[{{p{\linewidth}}}]{
                \texttt{\tiny{[GM$|$GM]}}
                \texttt{kraj} \\
            }
        }
    }
    & & \\ \\

    \theutterance \stepcounter{utterance}  
    & & & \multicolumn{2}{p{0.3\linewidth}}{
        \cellcolor[rgb]{0.9,0.9,0.9}{
            \makecell[{{p{\linewidth}}}]{
                \texttt{\tiny{[GM$|$GM]}}
                ['francja']
            }
        }
    }
    & & \\ \\

    \theutterance \stepcounter{utterance}  
    & & & \multicolumn{2}{p{0.3\linewidth}}{
        \cellcolor[rgb]{0.9,0.9,0.9}{
            \makecell[{{p{\linewidth}}}]{
                \texttt{\tiny{[GM$|$GM]}}
                {'word': 'francja', 'assignment': 'team'}
            }
        }
    }
    & & \\ \\

    \theutterance \stepcounter{utterance}  
    & & \multicolumn{4}{p{0.6\linewidth}}{
        \cellcolor[rgb]{0.9,0.9,0.9}{
            \makecell[{{p{\linewidth}}}]{
                \texttt{\tiny{[GM$\rangle$P2]}}
                \texttt{Słowo bugle zostało odgadnięte prawidłowo. Twoja tura się tam skończyła.} \\
\texttt{Teraz podaj kolejną listę oddzieloną przecinkami, zawierającą co najmniej 1 i maksymalnie 1 słów z poniższej listy słów, które najlepiej odnoszą się lub są najbardziej związane ze słowem 'kraj'. Pamiętaj, aby rozpocząć odpowiedź od ZGADNIĘCIE: . Uwaga: niektóre słowa zostały usunięte z listy w porównaniu do poprzednich próśb.} \\
\\ 
\texttt{oliwka, mucha, australia, indyk, taniec, ośmiornica, ryba, miedź, laboratorium, choroba, podanie, cykl, zegarek, czekolada, port, francja, dyrygent, dobrze, kłoda, twarz} \\
            }
        }
    }
    & \\ \\

    \theutterance \stepcounter{utterance}  
    & & & \multicolumn{4}{p{0.6\linewidth}}{
        \cellcolor[rgb]{0.9,0.9,0.9}{
            \makecell[{{p{\linewidth}}}]{
                \texttt{\tiny{[GM$\langle$P2]}}
                \texttt{ZGADNIĘCIE: australia} \\
            }
        }
    }
     \\ \\

    \theutterance \stepcounter{utterance}  
    & & & \multicolumn{2}{p{0.3\linewidth}}{
        \cellcolor[rgb]{0.9,0.9,0.9}{
            \makecell[{{p{\linewidth}}}]{
                \texttt{\tiny{[GM$|$GM]}}
                ['australia']
            }
        }
    }
    & & \\ \\

    \theutterance \stepcounter{utterance}  
    & & & \multicolumn{2}{p{0.3\linewidth}}{
        \cellcolor[rgb]{0.9,0.9,0.9}{
            \makecell[{{p{\linewidth}}}]{
                \texttt{\tiny{[GM$|$GM]}}
                {'word': 'australia', 'assignment': 'innocent'}
            }
        }
    }
    & & \\ \\

    \theutterance \stepcounter{utterance}  
    & & & \multicolumn{2}{p{0.3\linewidth}}{
        \cellcolor[rgb]{0.9,0.9,0.9}{
            \makecell[{{p{\linewidth}}}]{
                \texttt{\tiny{[GM$|$GM]}}
                \texttt{australia} \\
            }
        }
    }
    & & \\ \\

    \theutterance \stepcounter{utterance}  
    & & & \multicolumn{2}{p{0.3\linewidth}}{
        \cellcolor[rgb]{0.9,0.9,0.9}{
            \makecell[{{p{\linewidth}}}]{
                \texttt{\tiny{[GM$|$GM]}}
                {'word': 'dobrze', 'assignment': 'opponent'}
            }
        }
    }
    & & \\ \\

    \theutterance \stepcounter{utterance}  
    & & & \multicolumn{2}{p{0.3\linewidth}}{
        \cellcolor[rgb]{0.9,0.9,0.9}{
            \makecell[{{p{\linewidth}}}]{
                \texttt{\tiny{[GM$|$GM]}}
                {'hidden': {'team': ['taniec', 'port', 'indyk', 'ryba', 'miedź', 'francja', 'laboratorium'], 'innocent': ['ośmiornica', 'twarz', 'czekolada', 'zegarek', 'kłoda'], 'opponent': ['mucha', 'dyrygent', 'oliwka', 'podanie', 'cykl'], 'assassin': ['choroba']}, 'revealed': {'team': {'team': ['pralka', 'bugle'], 'innocent': ['helikopter', 'australia'], 'opponent': [], 'assassin': []}, 'opponent': {'team': [], 'innocent': [], 'opponent': ['spadochron', 'talerz', 'dobrze'], 'assassin': []}}}
            }
        }
    }
    & & \\ \\

    \theutterance \stepcounter{utterance}  
    & & \multicolumn{4}{p{0.6\linewidth}}{
        \cellcolor[rgb]{0.9,0.9,0.9}{
            \makecell[{{p{\linewidth}}}]{
                \texttt{\tiny{[P1$\langle$GM]}}
                \texttt{Słowo australia zostało odgadnięte, ale to słowo niewinny. Tura twojego kolegi z drużyny się tam skończyła.} \\
\texttt{Teraz podaj kolejną wskazówkę odnoszącą się do niektórych z pozostałych słów twojego zespołu oraz listę powiązanych słów docelowych. Pamiętaj, aby rozpocząć wskazówkę od 'WSKAZÓWKA: ', wstawić nową linię i rozpocząć listę słów docelowych oddzielonych przecinkami od 'CELE: '. Uwaga: niektóre słowa zostały usunięte z list w porównaniu do poprzednich żądań.} \\
\\ 
\texttt{Twoje pozostałe słowa zespołu to: taniec, port, indyk, ryba, miedź, francja, laboratorium.} \\
\texttt{Pozostałe słowa twojego przeciwnika to: mucha, dyrygent, oliwka, dobrze, podanie, cykl.} \\
\texttt{Pozostałe słowa rozpraszające to: ośmiornica, twarz, czekolada, zegarek, kłoda.} \\
\texttt{Pozostałe słowa zabójcy to: choroba.} \\
            }
        }
    }
    & \\ \\

    \theutterance \stepcounter{utterance}  
    & \multicolumn{4}{p{0.6\linewidth}}{
        \cellcolor[rgb]{0.9,0.9,0.9}{
            \makecell[{{p{\linewidth}}}]{
                \texttt{\tiny{[P1$\rangle$GM]}}
                \texttt{WSKAZÓWKA: foka} \\
\texttt{CELE: port, ryba, taniec} \\
            }
        }
    }
    & & \\ \\

    \theutterance \stepcounter{utterance}  
    & & & \multicolumn{2}{p{0.3\linewidth}}{
        \cellcolor[rgb]{0.9,0.9,0.9}{
            \makecell[{{p{\linewidth}}}]{
                \texttt{\tiny{[GM$|$GM]}}
                \texttt{foka} \\
            }
        }
    }
    & & \\ \\

    \theutterance \stepcounter{utterance}  
    & & & \multicolumn{2}{p{0.3\linewidth}}{
        \cellcolor[rgb]{0.9,0.9,0.9}{
            \makecell[{{p{\linewidth}}}]{
                \texttt{\tiny{[GM$|$GM]}}
                ['port', 'ryba', 'taniec']
            }
        }
    }
    & & \\ \\

    \theutterance \stepcounter{utterance}  
    & & & \multicolumn{2}{p{0.3\linewidth}}{
        \cellcolor[rgb]{0.9,0.9,0.9}{
            \makecell[{{p{\linewidth}}}]{
                \texttt{\tiny{[GM$|$GM]}}
                {'word': 'port', 'assignment': 'team'}
            }
        }
    }
    & & \\ \\

    \theutterance \stepcounter{utterance}  
    & & & \multicolumn{2}{p{0.3\linewidth}}{
        \cellcolor[rgb]{0.9,0.9,0.9}{
            \makecell[{{p{\linewidth}}}]{
                \texttt{\tiny{[GM$|$GM]}}
                {'word': 'ryba', 'assignment': 'team'}
            }
        }
    }
    & & \\ \\

    \theutterance \stepcounter{utterance}  
    & & & \multicolumn{2}{p{0.3\linewidth}}{
        \cellcolor[rgb]{0.9,0.9,0.9}{
            \makecell[{{p{\linewidth}}}]{
                \texttt{\tiny{[GM$|$GM]}}
                {'word': 'taniec', 'assignment': 'team'}
            }
        }
    }
    & & \\ \\

    \theutterance \stepcounter{utterance}  
    & & \multicolumn{4}{p{0.6\linewidth}}{
        \cellcolor[rgb]{0.9,0.9,0.9}{
            \makecell[{{p{\linewidth}}}]{
                \texttt{\tiny{[GM$\rangle$P2]}}
                \texttt{Słowo australia zostało odgadnięte, ale to słowo niewinny. Twoja tura się tam skończyła.} \\
\texttt{Teraz podaj kolejną listę oddzieloną przecinkami, zawierającą co najmniej 1 i maksymalnie 3 słów z poniższej listy słów, które najlepiej odnoszą się lub są najbardziej związane ze słowem 'foka'. Pamiętaj, aby rozpocząć odpowiedź od ZGADNIĘCIE: . Uwaga: niektóre słowa zostały usunięte z listy w porównaniu do poprzednich próśb.} \\
\\ 
\texttt{oliwka, mucha, indyk, taniec, ośmiornica, ryba, miedź, laboratorium, choroba, podanie, cykl, zegarek, czekolada, port, francja, dyrygent, kłoda, twarz} \\
            }
        }
    }
    & \\ \\

    \theutterance \stepcounter{utterance}  
    & & & \multicolumn{4}{p{0.6\linewidth}}{
        \cellcolor[rgb]{0.9,0.9,0.9}{
            \makecell[{{p{\linewidth}}}]{
                \texttt{\tiny{[GM$\langle$P2]}}
                \texttt{ZGADNIĘCIE: ryba} \\
            }
        }
    }
     \\ \\

    \theutterance \stepcounter{utterance}  
    & & & \multicolumn{2}{p{0.3\linewidth}}{
        \cellcolor[rgb]{0.9,0.9,0.9}{
            \makecell[{{p{\linewidth}}}]{
                \texttt{\tiny{[GM$|$GM]}}
                ['ryba']
            }
        }
    }
    & & \\ \\

    \theutterance \stepcounter{utterance}  
    & & & \multicolumn{2}{p{0.3\linewidth}}{
        \cellcolor[rgb]{0.9,0.9,0.9}{
            \makecell[{{p{\linewidth}}}]{
                \texttt{\tiny{[GM$|$GM]}}
                {'word': 'ryba', 'assignment': 'team'}
            }
        }
    }
    & & \\ \\

    \theutterance \stepcounter{utterance}  
    & & & \multicolumn{2}{p{0.3\linewidth}}{
        \cellcolor[rgb]{0.9,0.9,0.9}{
            \makecell[{{p{\linewidth}}}]{
                \texttt{\tiny{[GM$|$GM]}}
                {'word': 'ryba', 'assignment': 'team'}
            }
        }
    }
    & & \\ \\

    \theutterance \stepcounter{utterance}  
    & & & \multicolumn{2}{p{0.3\linewidth}}{
        \cellcolor[rgb]{0.9,0.9,0.9}{
            \makecell[{{p{\linewidth}}}]{
                \texttt{\tiny{[GM$|$GM]}}
                {'word': 'oliwka', 'assignment': 'opponent'}
            }
        }
    }
    & & \\ \\

    \theutterance \stepcounter{utterance}  
    & & & \multicolumn{2}{p{0.3\linewidth}}{
        \cellcolor[rgb]{0.9,0.9,0.9}{
            \makecell[{{p{\linewidth}}}]{
                \texttt{\tiny{[GM$|$GM]}}
                {'hidden': {'team': ['taniec', 'port', 'indyk', 'miedź', 'francja', 'laboratorium'], 'innocent': ['ośmiornica', 'twarz', 'czekolada', 'zegarek', 'kłoda'], 'opponent': ['mucha', 'dyrygent', 'podanie', 'cykl'], 'assassin': ['choroba']}, 'revealed': {'team': {'team': ['pralka', 'bugle', 'ryba'], 'innocent': ['helikopter', 'australia'], 'opponent': [], 'assassin': []}, 'opponent': {'team': [], 'innocent': [], 'opponent': ['spadochron', 'talerz', 'dobrze', 'oliwka'], 'assassin': []}}}
            }
        }
    }
    & & \\ \\

    \theutterance \stepcounter{utterance}  
    & & \multicolumn{4}{p{0.6\linewidth}}{
        \cellcolor[rgb]{0.9,0.9,0.9}{
            \makecell[{{p{\linewidth}}}]{
                \texttt{\tiny{[P1$\langle$GM]}}
                \texttt{Słowo ryba zostało odgadnięte prawidłowo. Tura twojego kolegi z drużyny się tam skończyła.} \\
\texttt{Teraz podaj kolejną wskazówkę odnoszącą się do niektórych z pozostałych słów twojego zespołu oraz listę powiązanych słów docelowych. Pamiętaj, aby rozpocząć wskazówkę od 'WSKAZÓWKA: ', wstawić nową linię i rozpocząć listę słów docelowych oddzielonych przecinkami od 'CELE: '. Uwaga: niektóre słowa zostały usunięte z list w porównaniu do poprzednich żądań.} \\
\\ 
\texttt{Twoje pozostałe słowa zespołu to: taniec, port, indyk, miedź, francja, laboratorium.} \\
\texttt{Pozostałe słowa twojego przeciwnika to: mucha, dyrygent, oliwka, podanie, cykl.} \\
\texttt{Pozostałe słowa rozpraszające to: ośmiornica, twarz, czekolada, zegarek, kłoda.} \\
\texttt{Pozostałe słowa zabójcy to: choroba.} \\
            }
        }
    }
    & \\ \\

    \theutterance \stepcounter{utterance}  
    & \multicolumn{4}{p{0.6\linewidth}}{
        \cellcolor[rgb]{0.9,0.9,0.9}{
            \makecell[{{p{\linewidth}}}]{
                \texttt{\tiny{[P1$\rangle$GM]}}
                \texttt{WSKAZÓWKA: metal} \\
\texttt{CELE: miedź, port} \\
            }
        }
    }
    & & \\ \\

    \theutterance \stepcounter{utterance}  
    & & & \multicolumn{2}{p{0.3\linewidth}}{
        \cellcolor[rgb]{0.9,0.9,0.9}{
            \makecell[{{p{\linewidth}}}]{
                \texttt{\tiny{[GM$|$GM]}}
                \texttt{metal} \\
            }
        }
    }
    & & \\ \\

    \theutterance \stepcounter{utterance}  
    & & & \multicolumn{2}{p{0.3\linewidth}}{
        \cellcolor[rgb]{0.9,0.9,0.9}{
            \makecell[{{p{\linewidth}}}]{
                \texttt{\tiny{[GM$|$GM]}}
                ['miedź', 'port']
            }
        }
    }
    & & \\ \\

    \theutterance \stepcounter{utterance}  
    & & & \multicolumn{2}{p{0.3\linewidth}}{
        \cellcolor[rgb]{0.9,0.9,0.9}{
            \makecell[{{p{\linewidth}}}]{
                \texttt{\tiny{[GM$|$GM]}}
                {'word': 'miedź', 'assignment': 'team'}
            }
        }
    }
    & & \\ \\

    \theutterance \stepcounter{utterance}  
    & & & \multicolumn{2}{p{0.3\linewidth}}{
        \cellcolor[rgb]{0.9,0.9,0.9}{
            \makecell[{{p{\linewidth}}}]{
                \texttt{\tiny{[GM$|$GM]}}
                {'word': 'port', 'assignment': 'team'}
            }
        }
    }
    & & \\ \\

    \theutterance \stepcounter{utterance}  
    & & \multicolumn{4}{p{0.6\linewidth}}{
        \cellcolor[rgb]{0.9,0.9,0.9}{
            \makecell[{{p{\linewidth}}}]{
                \texttt{\tiny{[GM$\rangle$P2]}}
                \texttt{Słowo ryba zostało odgadnięte prawidłowo. Twoja tura się tam skończyła.} \\
\texttt{Teraz podaj kolejną listę oddzieloną przecinkami, zawierającą co najmniej 1 i maksymalnie 2 słów z poniższej listy słów, które najlepiej odnoszą się lub są najbardziej związane ze słowem 'metal'. Pamiętaj, aby rozpocząć odpowiedź od ZGADNIĘCIE: . Uwaga: niektóre słowa zostały usunięte z listy w porównaniu do poprzednich próśb.} \\
\\ 
\texttt{mucha, indyk, taniec, ośmiornica, miedź, laboratorium, choroba, podanie, cykl, zegarek, czekolada, port, francja, dyrygent, kłoda, twarz} \\
            }
        }
    }
    & \\ \\

    \theutterance \stepcounter{utterance}  
    & & & \multicolumn{4}{p{0.6\linewidth}}{
        \cellcolor[rgb]{0.9,0.9,0.9}{
            \makecell[{{p{\linewidth}}}]{
                \texttt{\tiny{[GM$\langle$P2]}}
                \texttt{ZGADNIĘCIE: miedź, zegarek} \\
            }
        }
    }
     \\ \\

    \theutterance \stepcounter{utterance}  
    & & & \multicolumn{2}{p{0.3\linewidth}}{
        \cellcolor[rgb]{0.9,0.9,0.9}{
            \makecell[{{p{\linewidth}}}]{
                \texttt{\tiny{[GM$|$GM]}}
                ['miedź', 'zegarek']
            }
        }
    }
    & & \\ \\

    \theutterance \stepcounter{utterance}  
    & & & \multicolumn{2}{p{0.3\linewidth}}{
        \cellcolor[rgb]{0.9,0.9,0.9}{
            \makecell[{{p{\linewidth}}}]{
                \texttt{\tiny{[GM$|$GM]}}
                {'word': 'miedź', 'assignment': 'team'}
            }
        }
    }
    & & \\ \\

    \theutterance \stepcounter{utterance}  
    & & & \multicolumn{2}{p{0.3\linewidth}}{
        \cellcolor[rgb]{0.9,0.9,0.9}{
            \makecell[{{p{\linewidth}}}]{
                \texttt{\tiny{[GM$|$GM]}}
                {'word': 'miedź', 'assignment': 'team'}
            }
        }
    }
    & & \\ \\

    \theutterance \stepcounter{utterance}  
    & & & \multicolumn{2}{p{0.3\linewidth}}{
        \cellcolor[rgb]{0.9,0.9,0.9}{
            \makecell[{{p{\linewidth}}}]{
                \texttt{\tiny{[GM$|$GM]}}
                {'word': 'zegarek', 'assignment': 'innocent'}
            }
        }
    }
    & & \\ \\

    \theutterance \stepcounter{utterance}  
    & & & \multicolumn{2}{p{0.3\linewidth}}{
        \cellcolor[rgb]{0.9,0.9,0.9}{
            \makecell[{{p{\linewidth}}}]{
                \texttt{\tiny{[GM$|$GM]}}
                \texttt{zegarek} \\
            }
        }
    }
    & & \\ \\

    \theutterance \stepcounter{utterance}  
    & & & \multicolumn{2}{p{0.3\linewidth}}{
        \cellcolor[rgb]{0.9,0.9,0.9}{
            \makecell[{{p{\linewidth}}}]{
                \texttt{\tiny{[GM$|$GM]}}
                {'word': 'dyrygent', 'assignment': 'opponent'}
            }
        }
    }
    & & \\ \\

    \theutterance \stepcounter{utterance}  
    & & & \multicolumn{2}{p{0.3\linewidth}}{
        \cellcolor[rgb]{0.9,0.9,0.9}{
            \makecell[{{p{\linewidth}}}]{
                \texttt{\tiny{[GM$|$GM]}}
                {'hidden': {'team': ['taniec', 'port', 'indyk', 'francja', 'laboratorium'], 'innocent': ['ośmiornica', 'twarz', 'czekolada', 'kłoda'], 'opponent': ['mucha', 'podanie', 'cykl'], 'assassin': ['choroba']}, 'revealed': {'team': {'team': ['pralka', 'bugle', 'ryba', 'miedź'], 'innocent': ['helikopter', 'australia', 'zegarek'], 'opponent': [], 'assassin': []}, 'opponent': {'team': [], 'innocent': [], 'opponent': ['spadochron', 'talerz', 'dobrze', 'oliwka', 'dyrygent'], 'assassin': []}}}
            }
        }
    }
    & & \\ \\

    \theutterance \stepcounter{utterance}  
    & & \multicolumn{4}{p{0.6\linewidth}}{
        \cellcolor[rgb]{0.9,0.9,0.9}{
            \makecell[{{p{\linewidth}}}]{
                \texttt{\tiny{[P1$\langle$GM]}}
                \texttt{Słowo miedź zostało odgadnięte prawidłowo. Słowo zegarek zostało odgadnięte, ale to słowo niewinny. Tura twojego kolegi z drużyny się tam skończyła.} \\
\texttt{Teraz podaj kolejną wskazówkę odnoszącą się do niektórych z pozostałych słów twojego zespołu oraz listę powiązanych słów docelowych. Pamiętaj, aby rozpocząć wskazówkę od 'WSKAZÓWKA: ', wstawić nową linię i rozpocząć listę słów docelowych oddzielonych przecinkami od 'CELE: '. Uwaga: niektóre słowa zostały usunięte z list w porównaniu do poprzednich żądań.} \\
\\ 
\texttt{Twoje pozostałe słowa zespołu to: taniec, port, indyk, francja, laboratorium.} \\
\texttt{Pozostałe słowa twojego przeciwnika to: mucha, dyrygent, podanie, cykl.} \\
\texttt{Pozostałe słowa rozpraszające to: ośmiornica, twarz, czekolada, kłoda.} \\
\texttt{Pozostałe słowa zabójcy to: choroba.} \\
            }
        }
    }
    & \\ \\

    \theutterance \stepcounter{utterance}  
    & \multicolumn{4}{p{0.6\linewidth}}{
        \cellcolor[rgb]{0.9,0.9,0.9}{
            \makecell[{{p{\linewidth}}}]{
                \texttt{\tiny{[P1$\rangle$GM]}}
                \texttt{WSKAZÓWKA: kurort} \\
\texttt{CELE: francja, port} \\
            }
        }
    }
    & & \\ \\

    \theutterance \stepcounter{utterance}  
    & & & \multicolumn{2}{p{0.3\linewidth}}{
        \cellcolor[rgb]{0.9,0.9,0.9}{
            \makecell[{{p{\linewidth}}}]{
                \texttt{\tiny{[GM$|$GM]}}
                \texttt{kurort} \\
            }
        }
    }
    & & \\ \\

    \theutterance \stepcounter{utterance}  
    & & & \multicolumn{2}{p{0.3\linewidth}}{
        \cellcolor[rgb]{0.9,0.9,0.9}{
            \makecell[{{p{\linewidth}}}]{
                \texttt{\tiny{[GM$|$GM]}}
                ['francja', 'port']
            }
        }
    }
    & & \\ \\

    \theutterance \stepcounter{utterance}  
    & & & \multicolumn{2}{p{0.3\linewidth}}{
        \cellcolor[rgb]{0.9,0.9,0.9}{
            \makecell[{{p{\linewidth}}}]{
                \texttt{\tiny{[GM$|$GM]}}
                {'word': 'francja', 'assignment': 'team'}
            }
        }
    }
    & & \\ \\

    \theutterance \stepcounter{utterance}  
    & & & \multicolumn{2}{p{0.3\linewidth}}{
        \cellcolor[rgb]{0.9,0.9,0.9}{
            \makecell[{{p{\linewidth}}}]{
                \texttt{\tiny{[GM$|$GM]}}
                {'word': 'port', 'assignment': 'team'}
            }
        }
    }
    & & \\ \\

    \theutterance \stepcounter{utterance}  
    & & \multicolumn{4}{p{0.6\linewidth}}{
        \cellcolor[rgb]{0.9,0.9,0.9}{
            \makecell[{{p{\linewidth}}}]{
                \texttt{\tiny{[GM$\rangle$P2]}}
                \texttt{Słowo miedź zostało odgadnięte prawidłowo. Słowo zegarek zostało odgadnięte, ale to słowo niewinny. Twoja tura się tam skończyła.} \\
\texttt{Teraz podaj kolejną listę oddzieloną przecinkami, zawierającą co najmniej 1 i maksymalnie 2 słów z poniższej listy słów, które najlepiej odnoszą się lub są najbardziej związane ze słowem 'kurort'. Pamiętaj, aby rozpocząć odpowiedź od ZGADNIĘCIE: . Uwaga: niektóre słowa zostały usunięte z listy w porównaniu do poprzednich próśb.} \\
\\ 
\texttt{mucha, indyk, taniec, ośmiornica, laboratorium, choroba, podanie, cykl, czekolada, port, francja, kłoda, twarz} \\
            }
        }
    }
    & \\ \\

    \theutterance \stepcounter{utterance}  
    & & & \multicolumn{4}{p{0.6\linewidth}}{
        \cellcolor[rgb]{0.9,0.9,0.9}{
            \makecell[{{p{\linewidth}}}]{
                \texttt{\tiny{[GM$\langle$P2]}}
                \texttt{ZGADNIĘCIE: port, czekolada} \\
            }
        }
    }
     \\ \\

    \theutterance \stepcounter{utterance}  
    & & & \multicolumn{2}{p{0.3\linewidth}}{
        \cellcolor[rgb]{0.9,0.9,0.9}{
            \makecell[{{p{\linewidth}}}]{
                \texttt{\tiny{[GM$|$GM]}}
                ['port', 'czekolada']
            }
        }
    }
    & & \\ \\

    \theutterance \stepcounter{utterance}  
    & & & \multicolumn{2}{p{0.3\linewidth}}{
        \cellcolor[rgb]{0.9,0.9,0.9}{
            \makecell[{{p{\linewidth}}}]{
                \texttt{\tiny{[GM$|$GM]}}
                {'word': 'port', 'assignment': 'team'}
            }
        }
    }
    & & \\ \\

    \theutterance \stepcounter{utterance}  
    & & & \multicolumn{2}{p{0.3\linewidth}}{
        \cellcolor[rgb]{0.9,0.9,0.9}{
            \makecell[{{p{\linewidth}}}]{
                \texttt{\tiny{[GM$|$GM]}}
                {'word': 'port', 'assignment': 'team'}
            }
        }
    }
    & & \\ \\

    \theutterance \stepcounter{utterance}  
    & & & \multicolumn{2}{p{0.3\linewidth}}{
        \cellcolor[rgb]{0.9,0.9,0.9}{
            \makecell[{{p{\linewidth}}}]{
                \texttt{\tiny{[GM$|$GM]}}
                {'word': 'czekolada', 'assignment': 'innocent'}
            }
        }
    }
    & & \\ \\

    \theutterance \stepcounter{utterance}  
    & & & \multicolumn{2}{p{0.3\linewidth}}{
        \cellcolor[rgb]{0.9,0.9,0.9}{
            \makecell[{{p{\linewidth}}}]{
                \texttt{\tiny{[GM$|$GM]}}
                \texttt{czekolada} \\
            }
        }
    }
    & & \\ \\

    \theutterance \stepcounter{utterance}  
    & & & \multicolumn{2}{p{0.3\linewidth}}{
        \cellcolor[rgb]{0.9,0.9,0.9}{
            \makecell[{{p{\linewidth}}}]{
                \texttt{\tiny{[GM$|$GM]}}
                {'word': 'podanie', 'assignment': 'opponent'}
            }
        }
    }
    & & \\ \\

    \theutterance \stepcounter{utterance}  
    & & & \multicolumn{2}{p{0.3\linewidth}}{
        \cellcolor[rgb]{0.9,0.9,0.9}{
            \makecell[{{p{\linewidth}}}]{
                \texttt{\tiny{[GM$|$GM]}}
                {'hidden': {'team': ['taniec', 'indyk', 'francja', 'laboratorium'], 'innocent': ['ośmiornica', 'twarz', 'kłoda'], 'opponent': ['mucha', 'cykl'], 'assassin': ['choroba']}, 'revealed': {'team': {'team': ['pralka', 'bugle', 'ryba', 'miedź', 'port'], 'innocent': ['helikopter', 'australia', 'zegarek', 'czekolada'], 'opponent': [], 'assassin': []}, 'opponent': {'team': [], 'innocent': [], 'opponent': ['spadochron', 'talerz', 'dobrze', 'oliwka', 'dyrygent', 'podanie'], 'assassin': []}}}
            }
        }
    }
    & & \\ \\

    \theutterance \stepcounter{utterance}  
    & & \multicolumn{4}{p{0.6\linewidth}}{
        \cellcolor[rgb]{0.9,0.9,0.9}{
            \makecell[{{p{\linewidth}}}]{
                \texttt{\tiny{[P1$\langle$GM]}}
                \texttt{Słowo port zostało odgadnięte prawidłowo. Słowo czekolada zostało odgadnięte, ale to słowo niewinny. Tura twojego kolegi z drużyny się tam skończyła.} \\
\texttt{Teraz podaj kolejną wskazówkę odnoszącą się do niektórych z pozostałych słów twojego zespołu oraz listę powiązanych słów docelowych. Pamiętaj, aby rozpocząć wskazówkę od 'WSKAZÓWKA: ', wstawić nową linię i rozpocząć listę słów docelowych oddzielonych przecinkami od 'CELE: '. Uwaga: niektóre słowa zostały usunięte z list w porównaniu do poprzednich żądań.} \\
\\ 
\texttt{Twoje pozostałe słowa zespołu to: taniec, indyk, francja, laboratorium.} \\
\texttt{Pozostałe słowa twojego przeciwnika to: mucha, podanie, cykl.} \\
\texttt{Pozostałe słowa rozpraszające to: ośmiornica, twarz, kłoda.} \\
\texttt{Pozostałe słowa zabójcy to: choroba.} \\
            }
        }
    }
    & \\ \\

    \theutterance \stepcounter{utterance}  
    & \multicolumn{4}{p{0.6\linewidth}}{
        \cellcolor[rgb]{0.9,0.9,0.9}{
            \makecell[{{p{\linewidth}}}]{
                \texttt{\tiny{[P1$\rangle$GM]}}
                \texttt{WSKAZÓWKA: badania} \\
\texttt{CELE: laboratorium} \\
            }
        }
    }
    & & \\ \\

    \theutterance \stepcounter{utterance}  
    & & & \multicolumn{2}{p{0.3\linewidth}}{
        \cellcolor[rgb]{0.9,0.9,0.9}{
            \makecell[{{p{\linewidth}}}]{
                \texttt{\tiny{[GM$|$GM]}}
                \texttt{badania} \\
            }
        }
    }
    & & \\ \\

    \theutterance \stepcounter{utterance}  
    & & & \multicolumn{2}{p{0.3\linewidth}}{
        \cellcolor[rgb]{0.9,0.9,0.9}{
            \makecell[{{p{\linewidth}}}]{
                \texttt{\tiny{[GM$|$GM]}}
                ['laboratorium']
            }
        }
    }
    & & \\ \\

    \theutterance \stepcounter{utterance}  
    & & & \multicolumn{2}{p{0.3\linewidth}}{
        \cellcolor[rgb]{0.9,0.9,0.9}{
            \makecell[{{p{\linewidth}}}]{
                \texttt{\tiny{[GM$|$GM]}}
                {'word': 'laboratorium', 'assignment': 'team'}
            }
        }
    }
    & & \\ \\

    \theutterance \stepcounter{utterance}  
    & & \multicolumn{4}{p{0.6\linewidth}}{
        \cellcolor[rgb]{0.9,0.9,0.9}{
            \makecell[{{p{\linewidth}}}]{
                \texttt{\tiny{[GM$\rangle$P2]}}
                \texttt{Słowo port zostało odgadnięte prawidłowo. Słowo czekolada zostało odgadnięte, ale to słowo niewinny. Twoja tura się tam skończyła.} \\
\texttt{Teraz podaj kolejną listę oddzieloną przecinkami, zawierającą co najmniej 1 i maksymalnie 1 słów z poniższej listy słów, które najlepiej odnoszą się lub są najbardziej związane ze słowem 'badania'. Pamiętaj, aby rozpocząć odpowiedź od ZGADNIĘCIE: . Uwaga: niektóre słowa zostały usunięte z listy w porównaniu do poprzednich próśb.} \\
\\ 
\texttt{mucha, indyk, taniec, ośmiornica, laboratorium, choroba, cykl, francja, kłoda, twarz} \\
            }
        }
    }
    & \\ \\

    \theutterance \stepcounter{utterance}  
    & & & \multicolumn{4}{p{0.6\linewidth}}{
        \cellcolor[rgb]{0.9,0.9,0.9}{
            \makecell[{{p{\linewidth}}}]{
                \texttt{\tiny{[GM$\langle$P2]}}
                \texttt{ZGADNIĘCIE: laboratorium} \\
            }
        }
    }
     \\ \\

    \theutterance \stepcounter{utterance}  
    & & & \multicolumn{2}{p{0.3\linewidth}}{
        \cellcolor[rgb]{0.9,0.9,0.9}{
            \makecell[{{p{\linewidth}}}]{
                \texttt{\tiny{[GM$|$GM]}}
                ['laboratorium']
            }
        }
    }
    & & \\ \\

    \theutterance \stepcounter{utterance}  
    & & & \multicolumn{2}{p{0.3\linewidth}}{
        \cellcolor[rgb]{0.9,0.9,0.9}{
            \makecell[{{p{\linewidth}}}]{
                \texttt{\tiny{[GM$|$GM]}}
                {'word': 'laboratorium', 'assignment': 'team'}
            }
        }
    }
    & & \\ \\

    \theutterance \stepcounter{utterance}  
    & & & \multicolumn{2}{p{0.3\linewidth}}{
        \cellcolor[rgb]{0.9,0.9,0.9}{
            \makecell[{{p{\linewidth}}}]{
                \texttt{\tiny{[GM$|$GM]}}
                {'word': 'laboratorium', 'assignment': 'team'}
            }
        }
    }
    & & \\ \\

    \theutterance \stepcounter{utterance}  
    & & & \multicolumn{2}{p{0.3\linewidth}}{
        \cellcolor[rgb]{0.9,0.9,0.9}{
            \makecell[{{p{\linewidth}}}]{
                \texttt{\tiny{[GM$|$GM]}}
                {'word': 'cykl', 'assignment': 'opponent'}
            }
        }
    }
    & & \\ \\

    \theutterance \stepcounter{utterance}  
    & & & \multicolumn{2}{p{0.3\linewidth}}{
        \cellcolor[rgb]{0.9,0.9,0.9}{
            \makecell[{{p{\linewidth}}}]{
                \texttt{\tiny{[GM$|$GM]}}
                {'hidden': {'team': ['taniec', 'indyk', 'francja'], 'innocent': ['ośmiornica', 'twarz', 'kłoda'], 'opponent': ['mucha'], 'assassin': ['choroba']}, 'revealed': {'team': {'team': ['pralka', 'bugle', 'ryba', 'miedź', 'port', 'laboratorium'], 'innocent': ['helikopter', 'australia', 'zegarek', 'czekolada'], 'opponent': [], 'assassin': []}, 'opponent': {'team': [], 'innocent': [], 'opponent': ['spadochron', 'talerz', 'dobrze', 'oliwka', 'dyrygent', 'podanie', 'cykl'], 'assassin': []}}}
            }
        }
    }
    & & \\ \\

    \theutterance \stepcounter{utterance}  
    & & \multicolumn{4}{p{0.6\linewidth}}{
        \cellcolor[rgb]{0.9,0.9,0.9}{
            \makecell[{{p{\linewidth}}}]{
                \texttt{\tiny{[P1$\langle$GM]}}
                \texttt{Słowo laboratorium zostało odgadnięte prawidłowo. Tura twojego kolegi z drużyny się tam skończyła.} \\
\texttt{Teraz podaj kolejną wskazówkę odnoszącą się do niektórych z pozostałych słów twojego zespołu oraz listę powiązanych słów docelowych. Pamiętaj, aby rozpocząć wskazówkę od 'WSKAZÓWKA: ', wstawić nową linię i rozpocząć listę słów docelowych oddzielonych przecinkami od 'CELE: '. Uwaga: niektóre słowa zostały usunięte z list w porównaniu do poprzednich żądań.} \\
\\ 
\texttt{Twoje pozostałe słowa zespołu to: taniec, indyk, francja.} \\
\texttt{Pozostałe słowa twojego przeciwnika to: mucha, cykl.} \\
\texttt{Pozostałe słowa rozpraszające to: ośmiornica, twarz, kłoda.} \\
\texttt{Pozostałe słowa zabójcy to: choroba.} \\
            }
        }
    }
    & \\ \\

    \theutterance \stepcounter{utterance}  
    & \multicolumn{4}{p{0.6\linewidth}}{
        \cellcolor[rgb]{0.9,0.9,0.9}{
            \makecell[{{p{\linewidth}}}]{
                \texttt{\tiny{[P1$\rangle$GM]}}
                \texttt{WSKAZÓWKA: turcja} \\
\texttt{CELE: indyk, taniec} \\
            }
        }
    }
    & & \\ \\

    \theutterance \stepcounter{utterance}  
    & & & \multicolumn{2}{p{0.3\linewidth}}{
        \cellcolor[rgb]{0.9,0.9,0.9}{
            \makecell[{{p{\linewidth}}}]{
                \texttt{\tiny{[GM$|$GM]}}
                \texttt{turcja} \\
            }
        }
    }
    & & \\ \\

    \theutterance \stepcounter{utterance}  
    & & & \multicolumn{2}{p{0.3\linewidth}}{
        \cellcolor[rgb]{0.9,0.9,0.9}{
            \makecell[{{p{\linewidth}}}]{
                \texttt{\tiny{[GM$|$GM]}}
                ['indyk', 'taniec']
            }
        }
    }
    & & \\ \\

    \theutterance \stepcounter{utterance}  
    & & & \multicolumn{2}{p{0.3\linewidth}}{
        \cellcolor[rgb]{0.9,0.9,0.9}{
            \makecell[{{p{\linewidth}}}]{
                \texttt{\tiny{[GM$|$GM]}}
                {'word': 'indyk', 'assignment': 'team'}
            }
        }
    }
    & & \\ \\

    \theutterance \stepcounter{utterance}  
    & & & \multicolumn{2}{p{0.3\linewidth}}{
        \cellcolor[rgb]{0.9,0.9,0.9}{
            \makecell[{{p{\linewidth}}}]{
                \texttt{\tiny{[GM$|$GM]}}
                {'word': 'taniec', 'assignment': 'team'}
            }
        }
    }
    & & \\ \\

    \theutterance \stepcounter{utterance}  
    & & \multicolumn{4}{p{0.6\linewidth}}{
        \cellcolor[rgb]{0.9,0.9,0.9}{
            \makecell[{{p{\linewidth}}}]{
                \texttt{\tiny{[GM$\rangle$P2]}}
                \texttt{Słowo laboratorium zostało odgadnięte prawidłowo. Twoja tura się tam skończyła.} \\
\texttt{Teraz podaj kolejną listę oddzieloną przecinkami, zawierającą co najmniej 1 i maksymalnie 2 słów z poniższej listy słów, które najlepiej odnoszą się lub są najbardziej związane ze słowem 'turcja'. Pamiętaj, aby rozpocząć odpowiedź od ZGADNIĘCIE: . Uwaga: niektóre słowa zostały usunięte z listy w porównaniu do poprzednich próśb.} \\
\\ 
\texttt{mucha, indyk, taniec, ośmiornica, choroba, francja, kłoda, twarz} \\
            }
        }
    }
    & \\ \\

    \theutterance \stepcounter{utterance}  
    & & & \multicolumn{4}{p{0.6\linewidth}}{
        \cellcolor[rgb]{0.9,0.9,0.9}{
            \makecell[{{p{\linewidth}}}]{
                \texttt{\tiny{[GM$\langle$P2]}}
                \texttt{ZGADNIĘCIE: indyk} \\
            }
        }
    }
     \\ \\

    \theutterance \stepcounter{utterance}  
    & & & \multicolumn{2}{p{0.3\linewidth}}{
        \cellcolor[rgb]{0.9,0.9,0.9}{
            \makecell[{{p{\linewidth}}}]{
                \texttt{\tiny{[GM$|$GM]}}
                ['indyk']
            }
        }
    }
    & & \\ \\

    \theutterance \stepcounter{utterance}  
    & & & \multicolumn{2}{p{0.3\linewidth}}{
        \cellcolor[rgb]{0.9,0.9,0.9}{
            \makecell[{{p{\linewidth}}}]{
                \texttt{\tiny{[GM$|$GM]}}
                {'word': 'indyk', 'assignment': 'team'}
            }
        }
    }
    & & \\ \\

    \theutterance \stepcounter{utterance}  
    & & & \multicolumn{2}{p{0.3\linewidth}}{
        \cellcolor[rgb]{0.9,0.9,0.9}{
            \makecell[{{p{\linewidth}}}]{
                \texttt{\tiny{[GM$|$GM]}}
                {'word': 'indyk', 'assignment': 'team'}
            }
        }
    }
    & & \\ \\

    \theutterance \stepcounter{utterance}  
    & & & \multicolumn{2}{p{0.3\linewidth}}{
        \cellcolor[rgb]{0.9,0.9,0.9}{
            \makecell[{{p{\linewidth}}}]{
                \texttt{\tiny{[GM$|$GM]}}
                {'word': 'mucha', 'assignment': 'opponent'}
            }
        }
    }
    & & \\ \\

    \theutterance \stepcounter{utterance}  
    & & & \multicolumn{2}{p{0.3\linewidth}}{
        \cellcolor[rgb]{0.9,0.9,0.9}{
            \makecell[{{p{\linewidth}}}]{
                \texttt{\tiny{[GM$|$GM]}}
                {'hidden': {'team': ['taniec', 'francja'], 'innocent': ['ośmiornica', 'twarz', 'kłoda'], 'opponent': [], 'assassin': ['choroba']}, 'revealed': {'team': {'team': ['pralka', 'bugle', 'ryba', 'miedź', 'port', 'laboratorium', 'indyk'], 'innocent': ['helikopter', 'australia', 'zegarek', 'czekolada'], 'opponent': [], 'assassin': []}, 'opponent': {'team': [], 'innocent': [], 'opponent': ['spadochron', 'talerz', 'dobrze', 'oliwka', 'dyrygent', 'podanie', 'cykl', 'mucha'], 'assassin': []}}}
            }
        }
    }
    & & \\ \\

    \theutterance \stepcounter{utterance}  
    & & \multicolumn{4}{p{0.6\linewidth}}{
        \cellcolor[rgb]{0.9,0.9,0.9}{
            \makecell[{{p{\linewidth}}}]{
                \texttt{\tiny{[P1$\langle$GM]}}
                \texttt{Słowo indyk zostało odgadnięte prawidłowo. Tura twojego kolegi z drużyny się tam skończyła.} \\
\texttt{Teraz podaj kolejną wskazówkę odnoszącą się do niektórych z pozostałych słów twojego zespołu oraz listę powiązanych słów docelowych. Pamiętaj, aby rozpocząć wskazówkę od 'WSKAZÓWKA: ', wstawić nową linię i rozpocząć listę słów docelowych oddzielonych przecinkami od 'CELE: '. Uwaga: niektóre słowa zostały usunięte z list w porównaniu do poprzednich żądań.} \\
\\ 
\texttt{Twoje pozostałe słowa zespołu to: taniec, francja.} \\
\texttt{Pozostałe słowa twojego przeciwnika to: mucha.} \\
\texttt{Pozostałe słowa rozpraszające to: ośmiornica, twarz, kłoda.} \\
\texttt{Pozostałe słowa zabójcy to: choroba.} \\
            }
        }
    }
    & \\ \\

    \theutterance \stepcounter{utterance}  
    & \multicolumn{4}{p{0.6\linewidth}}{
        \cellcolor[rgb]{0.9,0.9,0.9}{
            \makecell[{{p{\linewidth}}}]{
                \texttt{\tiny{[P1$\rangle$GM]}}
                \texttt{WSKAZÓWKA: paryż} \\
\texttt{CELE: francja, taniec} \\
            }
        }
    }
    & & \\ \\

    \theutterance \stepcounter{utterance}  
    & & & \multicolumn{2}{p{0.3\linewidth}}{
        \cellcolor[rgb]{0.9,0.9,0.9}{
            \makecell[{{p{\linewidth}}}]{
                \texttt{\tiny{[GM$|$GM]}}
                \texttt{paryż} \\
            }
        }
    }
    & & \\ \\

    \theutterance \stepcounter{utterance}  
    & & & \multicolumn{2}{p{0.3\linewidth}}{
        \cellcolor[rgb]{0.9,0.9,0.9}{
            \makecell[{{p{\linewidth}}}]{
                \texttt{\tiny{[GM$|$GM]}}
                ['francja', 'taniec']
            }
        }
    }
    & & \\ \\

    \theutterance \stepcounter{utterance}  
    & & & \multicolumn{2}{p{0.3\linewidth}}{
        \cellcolor[rgb]{0.9,0.9,0.9}{
            \makecell[{{p{\linewidth}}}]{
                \texttt{\tiny{[GM$|$GM]}}
                {'word': 'francja', 'assignment': 'team'}
            }
        }
    }
    & & \\ \\

    \theutterance \stepcounter{utterance}  
    & & & \multicolumn{2}{p{0.3\linewidth}}{
        \cellcolor[rgb]{0.9,0.9,0.9}{
            \makecell[{{p{\linewidth}}}]{
                \texttt{\tiny{[GM$|$GM]}}
                {'word': 'taniec', 'assignment': 'team'}
            }
        }
    }
    & & \\ \\

    \theutterance \stepcounter{utterance}  
    & & & \multicolumn{2}{p{0.3\linewidth}}{
        \cellcolor[rgb]{0.9,0.9,0.9}{
            \makecell[{{p{\linewidth}}}]{
                \texttt{\tiny{[GM$|$GM]}}
                \texttt{opponent has won} \\
            }
        }
    }
    & & \\ \\

\end{supertabular}
}

\end{document}
